\begin{problem}[Four-digit numbers from 1 to 8]
How many $4$-digit numbers can you make using the digits $1$ to $8$ if no digit is repeated?
\end{problem}

\begin{solution}
There are $8$ choices for the thousands digit, then $7$ choices for the hundreds digit, $6$ for the tens digit, and $5$ for the units digit.
\[
8 \times 7 \times 6 \times 5 = 1680
\]
So the number of such $4$-digit numbers is $\boxed{1680}$.
\end{solution}

\begin{takeaways}
\item When order matters and repetition is not allowed, use a permutation count.
\item The multiplication principle is often the fastest method for straightforward arrangement questions.
\end{takeaways}

\begin{problem}[Rearranging the letters of COOKED]
How many different arrangements of the letters in the word \texttt{COOKED} are possible?
\end{problem}

\begin{solution}
The word \texttt{COOKED} has $6$ letters in total, with:
\[
O \text{ repeated } 2 \text{ times}.
\]
If all letters were distinct, there would be $6!$ arrangements. We divide by $2!$ for the repeated $O$'s:
\[
\frac{6!}{2!} = \frac{720}{2} = 360.
\]
Hence the number of different arrangements is $\boxed{360}$.
\end{solution}

\begin{takeaways}
\item For repeated letters, start with the total factorial and divide by the factorial of each repeated group.
\item This is the standard multiset permutation formula.
\end{takeaways}

\begin{problem}[Choosing three coloured balls]
A bag contains six different coloured balls.

How many different ways can three balls be selected?
\end{problem}

\begin{solution}
Since the order of selection does not matter, we use combinations:
\[
\binom{6}{3} = \frac{6 \cdot 5 \cdot 4}{3 \cdot 2 \cdot 1} = 20.
\]
So the number of different selections is $\boxed{20}$.
\end{solution}

\begin{takeaways}
\item Use combinations when order does not matter.
\item The phrase ``select two balls'' usually signals a combination, not a permutation.
\end{takeaways}

\begin{problem}[Rearranging the letters of CONDOBOLIN]
In how many different ways can all the letters of the word \texttt{CONDOBOLIN} be arranged in a line?
\end{problem}

\begin{solution}
The word \texttt{CONDOBOLIN} has $10$ letters in total, with:
\[
O \text{ repeated } 3 \text{ times}, \qquad N \text{ repeated } 2 \text{ times}.
\]
If all letters were distinct, there would be $10!$ arrangements. We divide by $3!$ for the repeated $O$'s and by $2!$ for the repeated $N$'s:
\[
\frac{10!}{3!\,2!}=\frac{3628800}{12}=302400.
\]
Hence the number of different arrangements is $\boxed{302400}$.
\end{solution}

\begin{takeaways}
\item This is another multiset permutation problem: start with $n!$ and divide by repeated factorials.
\item Always count repeated letters carefully before substituting into the formula.
\end{takeaways}
